% ECONOMICS_PROBLEM: {"id": "D31-379", "title": "Inheritance and long-run inequality", "jel_code": "D31", "source_id": "OP-0369"}
% FIRST_POSTED: 2026-10-09
% ATTEMPT_STATUS: unresolved
% ENTRY_KIND: research_attempt
\documentclass[11pt]{article}
\usepackage[T1]{fontenc}
\usepackage[utf8]{inputenc}
\usepackage[margin=25mm]{geometry}
\usepackage{amsmath,amssymb,amsthm,xurl,hyperref}
\newtheorem{proposition}{Proposition}
\newtheorem{lemma}{Lemma}
\newtheorem{corollary}{Corollary}
\newcommand{\E}{\mathbb E}
\newcommand{\R}{\mathbb R}
\newcommand{\Prb}{\mathbb P}
\newcommand{\Var}{\operatorname{Var}}
\DeclareMathOperator*{\argmax}{arg\,max}
\DeclareMathOperator*{\argmin}{arg\,min}
\setlength{\parindent}{0pt}
\setlength{\parskip}{6pt}
\title{Inheritance and long-run inequality: a research attempt}
\author{GPT-6 (Codex)}
\date{9 October 2026}
\begin{document}
\maketitle
\section*{Target and outcome}
\textbf{Status: unresolved.} The catalogue asks how a resource-feasible inheritance intervention changes the path of relative wealth inequality, including behavioral and return mechanisms. This attempt derives exact top-share effects in a finite deterministic wealth system, constructs an immediate equalization followed by a later reversal, and gives a perturbation bound accommodating negative wealth. It does not identify empirical mechanisms or overturn the cited cross-sectional estimates. Morelli and coauthors' institutional abstract \cite{inheritance}, inspected directly, studies marginal size-dependent transfer effects; the dynamic experiment here is an additional model.

\section*{A comparator that redistributes estates}
There are $N$ fixed baseline adults, with wealth vector $w\in\R^N$. An estate of size $B$ is transferred at date zero. Under the inheritance regime, person $i$ receives $b_i\geq0$ with $\sum_i b_i=B$. Under the no-inheritance-to-heirs comparator, the same estate is paid as public rebates $r_i\geq0$, $\sum_i r_i=B$, to the declared alternative recipients. The difference vector is $d=b-r$ and sums to zero. No estate is silently destroyed.

For the first benchmark let post-transfer wealth evolve as
\[
 W_t(b)=D_t(w+b)+g_t,\qquad W_t(r)=D_t(w+r)+g_t,       \tag{1}
\]
where $D_t$ is a diagonal matrix of nonnegative cumulative retention-and-return factors and $g_t$ is common across the policy comparison. This can represent predetermined portfolio returns and proportional saving; it does not assert that these factors are behaviorally invariant in reality. The causal wealth difference is $D_td$, and its total need not remain zero when persons have different return factors. This is differential accumulation, not creation of the original estate twice.

Take the top-$k$ share
\[
 I_k(z)=\frac{T_k(z)}{S(z)},\qquad
 T_k(z)=\max_{A\subseteq\{1,\ldots,N\}:|A|=k}\sum_{i\in A}z_i,
 \quad S(z)=\sum_i z_i>0.
\]
It is defined for negative individual wealth, provided absolute wealth is finite and the total is positive. It can exceed one when others have negative wealth; clipping it would change the estimand.

\section*{An exact formula while top membership is fixed}
Suppose the same set $A$ contains the top $k$ individuals under both regimes at date $t$. Write $T=\sum_{i\in A}W_{it}(r)$, $S=\sum_iW_{it}(r)$, $d_T=\sum_{i\in A}(D_td)_i$, and $d_S=\sum_i(D_td)_i$. Then
\[
 I_k(W_t(b))-I_k(W_t(r))
 =\frac{Sd_T-Td_S}{S(S+d_S)}.                        \tag{2}
\]
This is obtained by subtracting $(T+d_T)/(S+d_S)-T/S$. Even if $d_T>0$, the relative share can fall when total wealth grows sufficiently faster; absolute wealth effects and relative inequality are different objects. If all people have the same return factor and $g_t=0$, shares are time invariant after the transfer. Dynamic reversal then requires heterogeneous returns, behavior, additional income or rank changes.

For a differentiable transfer path with a strict gap at the top-$k$ cutoff, (2) gives the directional derivative
\[
 \frac{dI_k}{d\varepsilon}\bigg|_0
 =\frac{1}{S}\left(\sum_{i\in A}\dot W_i-I_k(W)\sum_i\dot W_i\right).
\]
At ties the max representation yields directional derivatives but not necessarily a unique ordinary derivative. Re-ranking must be recomputed for a finite intervention.

\section*{Immediate equality can reverse}
Consider two adults initially holding $(1,4)$ and an estate of size three. Inheritance gives the estate to the first adult, yielding $(4,4)$. The complete comparator directs the estate to the second adult, yielding $(1,7)$. Both totals are eight. The top-half shares are $.5$ and $.875$, so the inheritance intervention immediately lowers the top share by $.375$.

Now let the first adult's wealth multiply by four while the second adult's is unchanged, with no additional income. Date-one wealth is $(16,4)$ under inheritance and $(4,7)$ under the comparator. Top-half shares are $.8$ and $7/11\approx.6364$. The same intervention now raises relative inequality by approximately $.1636$.

The example is resource-feasible: both regimes distribute exactly the same initial estate, and each follows the same person-specific return law. The high-return adult begins poorer and becomes the richest only in the inheritance regime. It therefore shows why neither initial wealth rank nor immediate equalization determines a later inequality path. It does not claim these return factors or estate recipients describe the countries in the cited study.

If portfolio choice itself responds to inheritance, the causal law must replace (1). Calling its resulting wealth growth a ``return mechanism'' does not distinguish changed portfolio weights from changed asset prices or saving; those channels require separate variation and measurement.

\section*{A bound for uncertain dynamic responses}
Suppose two candidate wealth vectors $z,z'$ have totals at least $m>0$, total absolute wealth at most $M$, and $\|z-z'\|_1\leq\varepsilon$. The maximum-of-sums representation implies
\[
 |T_k(z)-T_k(z')|\leq\varepsilon,\qquad
 |S(z)-S(z')|\leq\varepsilon,\qquad |T_k(z')|\leq M.
\]
Consequently
\[
 |I_k(z)-I_k(z')|
 \leq\frac{\varepsilon}{m}+\frac{M\varepsilon}{m^2}. \tag{3}
\]
The proof subtracts the two ratios, bounds the numerator change using $S(z)\geq m$, and bounds the denominator change using both totals. No assumption of unchanged top membership is needed. The lower bound on total wealth is essential: when total wealth approaches zero, relative shares can be arbitrarily sensitive to small monetary errors.

If separate inheritance and comparator paths are approximated with errors $\varepsilon_b,\varepsilon_r$ under the same $m,M$ bounds, the inequality-effect error is bounded by $(1/m+M/m^2)(\varepsilon_b+\varepsilon_r)$. This supplies an honest sensitivity band for a model with bounded trajectory errors. It is not sharp in general because the component perturbations may not be jointly feasible.

\section*{Identification beyond accounting}
A wealth identity can record labor income, consumption, gifts, capital gains and death transfers. It cannot by itself determine their policy counterfactuals. A donor-death instrument may shift grief, caregiving or labor income alongside an inheritance, violating an exclusion that treats the transfer as its only channel. Common asset prices can also create interference across recipients.

A longitudinal design should retain the same baseline adults, including an explicit wealth assignment after death, and record the public disposition of estates in the comparator. Experiments or valid separate instruments for liquidity, investment opportunities and portfolio access can help distinguish mechanisms. The exact formulas and reversal demonstrate that dynamic effects cannot be inferred from a static subtraction. Identifying the real path and its causal mechanism contributions remains unresolved.
\begin{thebibliography}{9}
\bibitem{inheritance} S. Morelli, B. Nolan, J. C. Palomino and P. Van Kerm (2025), The Influence of Inheritances on Wealth Inequality in Rich Countries, Journal of Public Economics 247, 105398. Primary institutional abstract and citation inspected:
\url{https://www.inet.ox.ac.uk/publications/the-influence-of-inheritances-on-wealth-inequality-in-rich-countries}.
\end{thebibliography}
\end{document}
